% --- Archon physics formalization source begin ---
% archon:physics
% archon:covers PhyXMiniProblems/problem_phyx_mini_0535.lean
% archon:source-report reports/phyx_mini/problem_phyx_mini_0535.source.json

\chapter{Physics problem phyx\_mini\_0535}
\label{ch:PhyXMiniProblems_problem_phyx_mini_0535}

\paragraph{Problem source.}
\#\# Physical scenario

Consider a "crystal"consisting of two fixed ions of charge \textbackslash{}( +e \textbackslash{}) and two electrons as shown in figure.

\#\# Question

Find the value of \textbackslash{}( d \textbackslash{}) for which the total energy is a minimum.

\#\# Answer choices

- A: A: 62.1 pm

- B: B: 54.4 pm

- C: C: 21.8 pm

- D: D: 49.9 pm

\#\# Figure caption (auxiliary; use the image as primary evidence)

The image contains alternating blue and red spheres aligned horizontally. Each pair of adjacent spheres is connected by an arrow labeled "d," indicating distance or separation. Vertical black lines pass through each sphere, with the lines for the blue spheres extending upwards and the lines for the red spheres extending downwards. The arrows between the spheres show alternating directions.

\#\# Dataset metadata

Quantum Phenomena; Physical Model Grounding Reasoning, Numerical Reasoning

\paragraph{Recorded answer/context.}
D: D: 49.9 pm

\paragraph{Figure/image path.}
/root/proposal\_for\_physic/science-mango/phyx\_mini\_run/phyx\_data/test\_image/535.png

\paragraph{Formalization target.}
create a compiling Lean file with sorry bodies at `PhyXMiniProblems/problem\_phyx\_mini\_0535.lean`. The Lean declarations must preserve the physical quantities, dimensions or dimensional roles, figure labels, governing-law hypotheses, and final relation expressed by this problem.
Use Mathlib/Physlib names found through LeanExplore where available. If a physics API is missing, introduce faithful local abstractions rather than scalar placeholder aliases.

\begin{theorem}[Physics formalization target]
\label{thm:physics:phyx_mini_0535:target}
\lean{PhyXMiniProblems.ProblemPhyXMini0535.problem_phyx_mini_0535}
\uses{def:physics:phyx-mini-0535:phyxminiproblems-problemphyxmini0535-lengthquantity, def:physics:phyx-mini-0535:phyxminiproblems-problemphyxmini0535-lengthinmeters, def:physics:phyx-mini-0535:phyxminiproblems-problemphyxmini0535-lengthinpicometers, def:physics:phyx-mini-0535:phyxminiproblems-problemphyxmini0535-fourparticlecrystalsetup, def:physics:phyx-mini-0535:phyxminiproblems-problemphyxmini0535-matchestwoiontwoelectronscenario, def:physics:phyx-mini-0535:phyxminiproblems-problemphyxmini0535-matchesprimarycrystalfigure, def:physics:phyx-mini-0535:phyxminiproblems-problemphyxmini0535-usesreferencecrystalconstants, def:physics:phyx-mini-0535:phyxminiproblems-problemphyxmini0535-hasphysicalcrystalparameters, def:physics:phyx-mini-0535:phyxminiproblems-problemphyxmini0535-satisfiesrigidboxgroundstatelaw, def:physics:phyx-mini-0535:phyxminiproblems-problemphyxmini0535-satisfiescoulombandtotalenergylaws, def:physics:phyx-mini-0535:phyxminiproblems-problemphyxmini0535-isuniqueenergyminimizingseparation, def:physics:phyx-mini-0535:phyxminiproblems-problemphyxmini0535-recordeddatasetanswer, def:physics:phyx-mini-0535:phyxminiproblems-problemphyxmini0535-isuniqueclosestdisplayedseparation, def:physics:phyx-mini-0535:phyxminiproblems-problemphyxmini0535-stationaryseparationinmeters}
The assigned autoformalize agent should translate this physics problem into Lean declarations in the covered file, with theorem and lemma proof bodies written as `by sorry`.
\end{theorem}
\begin{proof}
This is an autoformalization task, not a proof task. Produce statements that can later be proved by the physics prover without weakening the physical model.
\end{proof}

% --- Archon declaration topology begin ---
\section{Declaration topology}
\begin{definition}[Length Quantity]
  \label{def:physics:phyx-mini-0535:phyxminiproblems-problemphyxmini0535-lengthquantity}
  \lean{PhyXMiniProblems.ProblemPhyXMini0535.LengthQuantity}
  A nonnegative, unit-independent physical length
\end{definition}
\begin{proof}
  This is the declared model definition of Length Quantity.
\end{proof}
\begin{definition}[Signed Length Quantity]
  \label{def:physics:phyx-mini-0535:phyxminiproblems-problemphyxmini0535-signedlengthquantity}
  \lean{PhyXMiniProblems.ProblemPhyXMini0535.SignedLengthQuantity}
  A signed, unit-independent one-dimensional position
\end{definition}
\begin{proof}
  This is the declared model definition of Signed Length Quantity.
\end{proof}
\begin{definition}[Mass Quantity]
  \label{def:physics:phyx-mini-0535:phyxminiproblems-problemphyxmini0535-massquantity}
  \lean{PhyXMiniProblems.ProblemPhyXMini0535.MassQuantity}
  A nonnegative physical mass
\end{definition}
\begin{proof}
  This is the declared model definition of Mass Quantity.
\end{proof}
\begin{definition}[Charge Quantity]
  \label{def:physics:phyx-mini-0535:phyxminiproblems-problemphyxmini0535-chargequantity}
  \lean{PhyXMiniProblems.ProblemPhyXMini0535.ChargeQuantity}
  A signed physical electric charge
\end{definition}
\begin{proof}
  This is the declared model definition of Charge Quantity.
\end{proof}
\begin{definition}[action Dimension]
  \label{def:physics:phyx-mini-0535:phyxminiproblems-problemphyxmini0535-actiondimension}
  \lean{PhyXMiniProblems.ProblemPhyXMini0535.actionDimension}
  The physical dimension M L² T⁻¹ of action
\end{definition}
\begin{proof}
  This is the declared model definition of action Dimension.
\end{proof}
\begin{definition}[Action Quantity]
  \label{def:physics:phyx-mini-0535:phyxminiproblems-problemphyxmini0535-actionquantity}
  \lean{PhyXMiniProblems.ProblemPhyXMini0535.ActionQuantity}
  \uses{def:physics:phyx-mini-0535:phyxminiproblems-problemphyxmini0535-actiondimension}
  A nonnegative physical action, used for the reduced Planck constant
\end{definition}
\begin{proof}
  This is the declared model definition of Action Quantity.
\end{proof}
\begin{definition}[length Readout]
  \label{def:physics:phyx-mini-0535:phyxminiproblems-problemphyxmini0535-lengthreadout}
  \lean{PhyXMiniProblems.ProblemPhyXMini0535.lengthReadout}
  \uses{def:physics:phyx-mini-0535:phyxminiproblems-problemphyxmini0535-lengthquantity}
  Read a nonnegative physical length in a selected Physlib length unit
\end{definition}
\begin{proof}
  This is the declared model definition of length Readout.
\end{proof}
\begin{definition}[signed Length Readout]
  \label{def:physics:phyx-mini-0535:phyxminiproblems-problemphyxmini0535-signedlengthreadout}
  \lean{PhyXMiniProblems.ProblemPhyXMini0535.signedLengthReadout}
  \uses{def:physics:phyx-mini-0535:phyxminiproblems-problemphyxmini0535-signedlengthquantity}
  Read a signed one-dimensional position in a selected length unit
\end{definition}
\begin{proof}
  This is the declared model definition of signed Length Readout.
\end{proof}
\begin{definition}[length In Meters]
  \label{def:physics:phyx-mini-0535:phyxminiproblems-problemphyxmini0535-lengthinmeters}
  \lean{PhyXMiniProblems.ProblemPhyXMini0535.lengthInMeters}
  \uses{def:physics:phyx-mini-0535:phyxminiproblems-problemphyxmini0535-lengthquantity, def:physics:phyx-mini-0535:phyxminiproblems-problemphyxmini0535-lengthreadout}
  SI metre readout of a physical length
\end{definition}
\begin{proof}
  This is the declared model definition of length In Meters.
\end{proof}
\begin{definition}[length In Picometers]
  \label{def:physics:phyx-mini-0535:phyxminiproblems-problemphyxmini0535-lengthinpicometers}
  \lean{PhyXMiniProblems.ProblemPhyXMini0535.lengthInPicometers}
  \uses{def:physics:phyx-mini-0535:phyxminiproblems-problemphyxmini0535-lengthquantity, def:physics:phyx-mini-0535:phyxminiproblems-problemphyxmini0535-lengthreadout}
  Picometre readout used by the displayed answers
\end{definition}
\begin{proof}
  This is the declared model definition of length In Picometers.
\end{proof}
\begin{definition}[mass In Kilograms]
  \label{def:physics:phyx-mini-0535:phyxminiproblems-problemphyxmini0535-massinkilograms}
  \lean{PhyXMiniProblems.ProblemPhyXMini0535.massInKilograms}
  \uses{def:physics:phyx-mini-0535:phyxminiproblems-problemphyxmini0535-massquantity}
  SI kilogram readout of a physical mass
\end{definition}
\begin{proof}
  This is the declared model definition of mass In Kilograms.
\end{proof}
\begin{definition}[charge Readout]
  \label{def:physics:phyx-mini-0535:phyxminiproblems-problemphyxmini0535-chargereadout}
  \lean{PhyXMiniProblems.ProblemPhyXMini0535.chargeReadout}
  \uses{def:physics:phyx-mini-0535:phyxminiproblems-problemphyxmini0535-chargequantity}
  Read a charge in a selected Physlib charge unit
\end{definition}
\begin{proof}
  This is the declared model definition of charge Readout.
\end{proof}
\begin{definition}[charge In Coulombs]
  \label{def:physics:phyx-mini-0535:phyxminiproblems-problemphyxmini0535-chargeincoulombs}
  \lean{PhyXMiniProblems.ProblemPhyXMini0535.chargeInCoulombs}
  \uses{def:physics:phyx-mini-0535:phyxminiproblems-problemphyxmini0535-chargequantity, def:physics:phyx-mini-0535:phyxminiproblems-problemphyxmini0535-chargereadout}
  SI coulomb readout of a physical charge
\end{definition}
\begin{proof}
  This is the declared model definition of charge In Coulombs.
\end{proof}
\begin{definition}[charge In Elementary Charges]
  \label{def:physics:phyx-mini-0535:phyxminiproblems-problemphyxmini0535-chargeinelementarycharges}
  \lean{PhyXMiniProblems.ProblemPhyXMini0535.chargeInElementaryCharges}
  \uses{def:physics:phyx-mini-0535:phyxminiproblems-problemphyxmini0535-chargequantity, def:physics:phyx-mini-0535:phyxminiproblems-problemphyxmini0535-chargereadout}
  Read a charge as a signed multiple of the elementary charge
\end{definition}
\begin{proof}
  This is the declared model definition of charge In Elementary Charges.
\end{proof}
\begin{definition}[action In Joule Seconds]
  \label{def:physics:phyx-mini-0535:phyxminiproblems-problemphyxmini0535-actioninjouleseconds}
  \lean{PhyXMiniProblems.ProblemPhyXMini0535.actionInJouleSeconds}
  \uses{def:physics:phyx-mini-0535:phyxminiproblems-problemphyxmini0535-actionquantity}
  Coherent-SI readout of an action, in joule-seconds
\end{definition}
\begin{proof}
  This is the declared model definition of action In Joule Seconds.
\end{proof}
\begin{definition}[energy In Joules]
  \label{def:physics:phyx-mini-0535:phyxminiproblems-problemphyxmini0535-energyinjoules}
  \lean{PhyXMiniProblems.ProblemPhyXMini0535.energyInJoules}
  Coherent-SI readout of a Physlib energy, in joules
\end{definition}
\begin{proof}
  This is the declared model definition of energy In Joules.
\end{proof}
\begin{definition}[Crystal Particle]
  \label{def:physics:phyx-mini-0535:phyxminiproblems-problemphyxmini0535-crystalparticle}
  \lean{PhyXMiniProblems.ProblemPhyXMini0535.CrystalParticle}
  The four particles, named in their left-to-right order in the image
\end{definition}
\begin{proof}
  This is the declared model definition of Crystal Particle.
\end{proof}
\begin{definition}[Electron Label]
  \label{def:physics:phyx-mini-0535:phyxminiproblems-problemphyxmini0535-electronlabel}
  \lean{PhyXMiniProblems.ProblemPhyXMini0535.ElectronLabel}
  The two indistinguishable electrons, distinguished only by image position
\end{definition}
\begin{proof}
  This is the declared model definition of Electron Label.
\end{proof}
\begin{definition}[particle]
  \label{def:physics:phyx-mini-0535:phyxminiproblems-problemphyxmini0535-electronlabel-particle}
  \lean{PhyXMiniProblems.ProblemPhyXMini0535.ElectronLabel.particle}
  \uses{def:physics:phyx-mini-0535:phyxminiproblems-problemphyxmini0535-crystalparticle, def:physics:phyx-mini-0535:phyxminiproblems-problemphyxmini0535-electronlabel}
  Regard a labelled electron as one of the four pictured particles
\end{definition}
\begin{proof}
  This is the declared model definition of particle.
\end{proof}
\begin{definition}[Particle Species]
  \label{def:physics:phyx-mini-0535:phyxminiproblems-problemphyxmini0535-particlespecies}
  \lean{PhyXMiniProblems.ProblemPhyXMini0535.ParticleSpecies}
  The particle species stated in the prose
\end{definition}
\begin{proof}
  This is the declared model definition of Particle Species.
\end{proof}
\begin{definition}[Electron Spin]
  \label{def:physics:phyx-mini-0535:phyxminiproblems-problemphyxmini0535-electronspin}
  \lean{PhyXMiniProblems.ProblemPhyXMini0535.ElectronSpin}
  Opposite spin states allow the two electrons to share one spatial mode
\end{definition}
\begin{proof}
  This is the declared model definition of Electron Spin.
\end{proof}
\begin{definition}[Confinement Model]
  \label{def:physics:phyx-mini-0535:phyxminiproblems-problemphyxmini0535-confinementmodel}
  \lean{PhyXMiniProblems.ProblemPhyXMini0535.ConfinementModel}
  The confinement idealization used for the quantum kinetic energy
\end{definition}
\begin{proof}
  This is the declared model definition of Confinement Model.
\end{proof}
\begin{definition}[Particle Pair]
  \label{def:physics:phyx-mini-0535:phyxminiproblems-problemphyxmini0535-particlepair}
  \lean{PhyXMiniProblems.ProblemPhyXMini0535.ParticlePair}
  The six unordered pairs among the four pictured particles
\end{definition}
\begin{proof}
  This is the declared model definition of Particle Pair.
\end{proof}
\begin{definition}[first]
  \label{def:physics:phyx-mini-0535:phyxminiproblems-problemphyxmini0535-particlepair-first}
  \lean{PhyXMiniProblems.ProblemPhyXMini0535.ParticlePair.first}
  \uses{def:physics:phyx-mini-0535:phyxminiproblems-problemphyxmini0535-crystalparticle, def:physics:phyx-mini-0535:phyxminiproblems-problemphyxmini0535-particlepair}
  First endpoint of each unordered particle pair
\end{definition}
\begin{proof}
  This is the declared model definition of first.
\end{proof}
\begin{definition}[second]
  \label{def:physics:phyx-mini-0535:phyxminiproblems-problemphyxmini0535-particlepair-second}
  \lean{PhyXMiniProblems.ProblemPhyXMini0535.ParticlePair.second}
  \uses{def:physics:phyx-mini-0535:phyxminiproblems-problemphyxmini0535-crystalparticle, def:physics:phyx-mini-0535:phyxminiproblems-problemphyxmini0535-particlepair}
  Second endpoint of each unordered particle pair
\end{definition}
\begin{proof}
  This is the declared model definition of second.
\end{proof}
\begin{definition}[gap Count]
  \label{def:physics:phyx-mini-0535:phyxminiproblems-problemphyxmini0535-particlepair-gapcount}
  \lean{PhyXMiniProblems.ProblemPhyXMini0535.ParticlePair.gapCount}
  \uses{def:physics:phyx-mini-0535:phyxminiproblems-problemphyxmini0535-particlepair}
  Number of adjacent d-gaps separating the endpoints of a pair
\end{definition}
\begin{proof}
  This is the declared model definition of gap Count.
\end{proof}
\begin{definition}[Adjacent Gap]
  \label{def:physics:phyx-mini-0535:phyxminiproblems-problemphyxmini0535-adjacentgap}
  \lean{PhyXMiniProblems.ProblemPhyXMini0535.AdjacentGap}
  The three adjacent gaps explicitly marked in the image
\end{definition}
\begin{proof}
  This is the declared model definition of Adjacent Gap.
\end{proof}
\begin{definition}[Sphere Style]
  \label{def:physics:phyx-mini-0535:phyxminiproblems-problemphyxmini0535-spherestyle}
  \lean{PhyXMiniProblems.ProblemPhyXMini0535.SphereStyle}
  Relative size and colour of a sphere in the supplied raster
\end{definition}
\begin{proof}
  This is the declared model definition of Sphere Style.
\end{proof}
\begin{definition}[Figure Label]
  \label{def:physics:phyx-mini-0535:phyxminiproblems-problemphyxmini0535-figurelabel}
  \lean{PhyXMiniProblems.ProblemPhyXMini0535.FigureLabel}
  The sole symbolic length label visible in the supplied raster
\end{definition}
\begin{proof}
  This is the declared model definition of Figure Label.
\end{proof}
\begin{definition}[Separation Marker]
  \label{def:physics:phyx-mini-0535:phyxminiproblems-problemphyxmini0535-separationmarker}
  \lean{PhyXMiniProblems.ProblemPhyXMini0535.SeparationMarker}
  The image marks each adjacent gap in both horizontal directions
\end{definition}
\begin{proof}
  This is the declared model definition of Separation Marker.
\end{proof}
\begin{definition}[Four Particle Crystal Figure]
  \label{def:physics:phyx-mini-0535:phyxminiproblems-problemphyxmini0535-fourparticlecrystalfigure}
  \lean{PhyXMiniProblems.ProblemPhyXMini0535.FourParticleCrystalFigure}
  \uses{def:physics:phyx-mini-0535:phyxminiproblems-problemphyxmini0535-crystalparticle, def:physics:phyx-mini-0535:phyxminiproblems-problemphyxmini0535-adjacentgap, def:physics:phyx-mini-0535:phyxminiproblems-problemphyxmini0535-spherestyle, def:physics:phyx-mini-0535:phyxminiproblems-problemphyxmini0535-figurelabel, def:physics:phyx-mini-0535:phyxminiproblems-problemphyxmini0535-separationmarker}
  Qualitative information transcribed from the primary bitmap
\end{definition}
\begin{proof}
  This is the declared model definition of Four Particle Crystal Figure.
\end{proof}
\begin{definition}[Four Particle Crystal Setup]
  \label{def:physics:phyx-mini-0535:phyxminiproblems-problemphyxmini0535-fourparticlecrystalsetup}
  \lean{PhyXMiniProblems.ProblemPhyXMini0535.FourParticleCrystalSetup}
  \uses{def:physics:phyx-mini-0535:phyxminiproblems-problemphyxmini0535-lengthquantity, def:physics:phyx-mini-0535:phyxminiproblems-problemphyxmini0535-signedlengthquantity, def:physics:phyx-mini-0535:phyxminiproblems-problemphyxmini0535-massquantity, def:physics:phyx-mini-0535:phyxminiproblems-problemphyxmini0535-chargequantity, def:physics:phyx-mini-0535:phyxminiproblems-problemphyxmini0535-actionquantity, def:physics:phyx-mini-0535:phyxminiproblems-problemphyxmini0535-crystalparticle, def:physics:phyx-mini-0535:phyxminiproblems-problemphyxmini0535-electronlabel, def:physics:phyx-mini-0535:phyxminiproblems-problemphyxmini0535-particlespecies, def:physics:phyx-mini-0535:phyxminiproblems-problemphyxmini0535-electronspin, def:physics:phyx-mini-0535:phyxminiproblems-problemphyxmini0535-confinementmodel, def:physics:phyx-mini-0535:phyxminiproblems-problemphyxmini0535-particlepair, def:physics:phyx-mini-0535:phyxminiproblems-problemphyxmini0535-fourparticlecrystalfigure}
  The physical state and the energy observables for every trial separation. None of these fields defines a minimizing distance or mentions an answer choice
\end{definition}
\begin{proof}
  This is the declared model definition of Four Particle Crystal Setup.
\end{proof}
\begin{definition}[Matches Two Ion Two Electron Scenario]
  \label{def:physics:phyx-mini-0535:phyxminiproblems-problemphyxmini0535-matchestwoiontwoelectronscenario}
  \lean{PhyXMiniProblems.ProblemPhyXMini0535.MatchesTwoIonTwoElectronScenario}
  \uses{def:physics:phyx-mini-0535:phyxminiproblems-problemphyxmini0535-chargeinelementarycharges, def:physics:phyx-mini-0535:phyxminiproblems-problemphyxmini0535-electronlabel, def:physics:phyx-mini-0535:phyxminiproblems-problemphyxmini0535-fourparticlecrystalsetup}
  The two-fixed-ion, two-electron scenario stated in the problem
\end{definition}
\begin{proof}
  This is the declared model definition of Matches Two Ion Two Electron Scenario.
\end{proof}
\begin{definition}[Matches Primary Crystal Figure]
  \label{def:physics:phyx-mini-0535:phyxminiproblems-problemphyxmini0535-matchesprimarycrystalfigure}
  \lean{PhyXMiniProblems.ProblemPhyXMini0535.MatchesPrimaryCrystalFigure}
  \uses{def:physics:phyx-mini-0535:phyxminiproblems-problemphyxmini0535-lengthquantity, def:physics:phyx-mini-0535:phyxminiproblems-problemphyxmini0535-lengthreadout, def:physics:phyx-mini-0535:phyxminiproblems-problemphyxmini0535-signedlengthreadout, def:physics:phyx-mini-0535:phyxminiproblems-problemphyxmini0535-particlepair, def:physics:phyx-mini-0535:phyxminiproblems-problemphyxmini0535-fourparticlecrystalsetup}
  Evidence from the primary image. It fixes the alternating styles, the three d labels, and the linear coordinates for an arbitrary trial d; it does not select any numerical value of d
\end{definition}
\begin{proof}
  This is the declared model definition of Matches Primary Crystal Figure.
\end{proof}
\begin{definition}[Uses Reference Crystal Constants]
  \label{def:physics:phyx-mini-0535:phyxminiproblems-problemphyxmini0535-usesreferencecrystalconstants}
  \lean{PhyXMiniProblems.ProblemPhyXMini0535.UsesReferenceCrystalConstants}
  \uses{def:physics:phyx-mini-0535:phyxminiproblems-problemphyxmini0535-massinkilograms, def:physics:phyx-mini-0535:phyxminiproblems-problemphyxmini0535-actioninjouleseconds, def:physics:phyx-mini-0535:phyxminiproblems-problemphyxmini0535-fourparticlecrystalsetup}
  Reference constants used for the numerical comparison. Physlib supplies Constants.ℏ, ChargeUnit.elementaryCharge, and Coulomb's constant as the method Electromagnetism.EMSystem.coulombConstant. The electron mass and the free-space electromagnetic-system calibration are explicit SI data
\end{definition}
\begin{proof}
  This is the declared model definition of Uses Reference Crystal Constants.
\end{proof}
\begin{definition}[Has Physical Crystal Parameters]
  \label{def:physics:phyx-mini-0535:phyxminiproblems-problemphyxmini0535-hasphysicalcrystalparameters}
  \lean{PhyXMiniProblems.ProblemPhyXMini0535.HasPhysicalCrystalParameters}
  \uses{def:physics:phyx-mini-0535:phyxminiproblems-problemphyxmini0535-lengthquantity, def:physics:phyx-mini-0535:phyxminiproblems-problemphyxmini0535-lengthinmeters, def:physics:phyx-mini-0535:phyxminiproblems-problemphyxmini0535-massinkilograms, def:physics:phyx-mini-0535:phyxminiproblems-problemphyxmini0535-chargeincoulombs, def:physics:phyx-mini-0535:phyxminiproblems-problemphyxmini0535-actioninjouleseconds, def:physics:phyx-mini-0535:phyxminiproblems-problemphyxmini0535-particlepair, def:physics:phyx-mini-0535:phyxminiproblems-problemphyxmini0535-fourparticlecrystalsetup}
  Positivity conditions selecting a nondegenerate physical model
\end{definition}
\begin{proof}
  This is the declared model definition of Has Physical Crystal Parameters.
\end{proof}
\begin{definition}[Satisfies Rigid Box Ground State Law]
  \label{def:physics:phyx-mini-0535:phyxminiproblems-problemphyxmini0535-satisfiesrigidboxgroundstatelaw}
  \lean{PhyXMiniProblems.ProblemPhyXMini0535.SatisfiesRigidBoxGroundStateLaw}
  \uses{def:physics:phyx-mini-0535:phyxminiproblems-problemphyxmini0535-lengthquantity, def:physics:phyx-mini-0535:phyxminiproblems-problemphyxmini0535-lengthinmeters, def:physics:phyx-mini-0535:phyxminiproblems-problemphyxmini0535-massinkilograms, def:physics:phyx-mini-0535:phyxminiproblems-problemphyxmini0535-actioninjouleseconds, def:physics:phyx-mini-0535:phyxminiproblems-problemphyxmini0535-energyinjoules, def:physics:phyx-mini-0535:phyxminiproblems-problemphyxmini0535-electronlabel, def:physics:phyx-mini-0535:phyxminiproblems-problemphyxmini0535-fourparticlecrystalsetup}
  Both opposite-spin electrons occupy the n = 1 spatial mode. In a one-dimensional infinite well of length L, each then has kinetic energy (pi*hbar)\textasciicircum{}2/(2*m*L\textasciicircum{}2). This all-separations law contains no minimizing distance
\end{definition}
\begin{proof}
  This is the declared model definition of Satisfies Rigid Box Ground State Law.
\end{proof}
\begin{definition}[Satisfies Coulomb And Total Energy Laws]
  \label{def:physics:phyx-mini-0535:phyxminiproblems-problemphyxmini0535-satisfiescoulombandtotalenergylaws}
  \lean{PhyXMiniProblems.ProblemPhyXMini0535.SatisfiesCoulombAndTotalEnergyLaws}
  \uses{def:physics:phyx-mini-0535:phyxminiproblems-problemphyxmini0535-lengthquantity, def:physics:phyx-mini-0535:phyxminiproblems-problemphyxmini0535-lengthinmeters, def:physics:phyx-mini-0535:phyxminiproblems-problemphyxmini0535-chargeincoulombs, def:physics:phyx-mini-0535:phyxminiproblems-problemphyxmini0535-energyinjoules, def:physics:phyx-mini-0535:phyxminiproblems-problemphyxmini0535-electronlabel, def:physics:phyx-mini-0535:phyxminiproblems-problemphyxmini0535-particlepair, def:physics:phyx-mini-0535:phyxminiproblems-problemphyxmini0535-fourparticlecrystalsetup}
  Every pair contributes k q₁ q₂/r, and the total energy is the sum of the two electron kinetic energies and all six pair energies. These are general governing laws over positive trial separations, not equilibrium conditions
\end{definition}
\begin{proof}
  This is the declared model definition of Satisfies Coulomb And Total Energy Laws.
\end{proof}
\begin{definition}[total Energy In Joules At Spacing]
  \label{def:physics:phyx-mini-0535:phyxminiproblems-problemphyxmini0535-totalenergyinjoulesatspacing}
  \lean{PhyXMiniProblems.ProblemPhyXMini0535.totalEnergyInJoulesAtSpacing}
  \uses{def:physics:phyx-mini-0535:phyxminiproblems-problemphyxmini0535-lengthquantity, def:physics:phyx-mini-0535:phyxminiproblems-problemphyxmini0535-energyinjoules, def:physics:phyx-mini-0535:phyxminiproblems-problemphyxmini0535-fourparticlecrystalsetup}
  The scalar total-energy curve obtained by reading a trial energy in joules
\end{definition}
\begin{proof}
  This is the declared model definition of total Energy In Joules At Spacing.
\end{proof}
\begin{definition}[Is Energy Minimizing Separation]
  \label{def:physics:phyx-mini-0535:phyxminiproblems-problemphyxmini0535-isenergyminimizingseparation}
  \lean{PhyXMiniProblems.ProblemPhyXMini0535.IsEnergyMinimizingSeparation}
  \uses{def:physics:phyx-mini-0535:phyxminiproblems-problemphyxmini0535-lengthquantity, def:physics:phyx-mini-0535:phyxminiproblems-problemphyxmini0535-lengthinmeters, def:physics:phyx-mini-0535:phyxminiproblems-problemphyxmini0535-fourparticlecrystalsetup, def:physics:phyx-mini-0535:phyxminiproblems-problemphyxmini0535-totalenergyinjoulesatspacing}
  A positive separation at which the energy is globally minimal
\end{definition}
\begin{proof}
  This is the declared model definition of Is Energy Minimizing Separation.
\end{proof}
\begin{definition}[Is Unique Energy Minimizing Separation]
  \label{def:physics:phyx-mini-0535:phyxminiproblems-problemphyxmini0535-isuniqueenergyminimizingseparation}
  \lean{PhyXMiniProblems.ProblemPhyXMini0535.IsUniqueEnergyMinimizingSeparation}
  \uses{def:physics:phyx-mini-0535:phyxminiproblems-problemphyxmini0535-lengthquantity, def:physics:phyx-mini-0535:phyxminiproblems-problemphyxmini0535-lengthinmeters, def:physics:phyx-mini-0535:phyxminiproblems-problemphyxmini0535-fourparticlecrystalsetup, def:physics:phyx-mini-0535:phyxminiproblems-problemphyxmini0535-isenergyminimizingseparation}
  Global minimality together with uniqueness of the physical length readout
\end{definition}
\begin{proof}
  This is the declared model definition of Is Unique Energy Minimizing Separation.
\end{proof}
\begin{definition}[Answer Choice]
  \label{def:physics:phyx-mini-0535:phyxminiproblems-problemphyxmini0535-answerchoice}
  \lean{PhyXMiniProblems.ProblemPhyXMini0535.AnswerChoice}
  Labels of the four separation choices printed with the source problem
\end{definition}
\begin{proof}
  This is the declared model definition of Answer Choice.
\end{proof}
\begin{definition}[separation In Picometers]
  \label{def:physics:phyx-mini-0535:phyxminiproblems-problemphyxmini0535-answerchoice-separationinpicometers}
  \lean{PhyXMiniProblems.ProblemPhyXMini0535.AnswerChoice.separationInPicometers}
  \uses{def:physics:phyx-mini-0535:phyxminiproblems-problemphyxmini0535-answerchoice}
  Picometre value printed beside each answer label
\end{definition}
\begin{proof}
  This is the declared model definition of separation In Picometers.
\end{proof}
\begin{definition}[recorded Dataset Answer]
  \label{def:physics:phyx-mini-0535:phyxminiproblems-problemphyxmini0535-recordeddatasetanswer}
  \lean{PhyXMiniProblems.ProblemPhyXMini0535.recordedDatasetAnswer}
  \uses{def:physics:phyx-mini-0535:phyxminiproblems-problemphyxmini0535-answerchoice}
  The answer label recorded by the source dataset
\end{definition}
\begin{proof}
  This is the declared model definition of recorded Dataset Answer.
\end{proof}
\begin{definition}[Is Unique Closest Displayed Separation]
  \label{def:physics:phyx-mini-0535:phyxminiproblems-problemphyxmini0535-isuniqueclosestdisplayedseparation}
  \lean{PhyXMiniProblems.ProblemPhyXMini0535.IsUniqueClosestDisplayedSeparation}
  \uses{def:physics:phyx-mini-0535:phyxminiproblems-problemphyxmini0535-lengthquantity, def:physics:phyx-mini-0535:phyxminiproblems-problemphyxmini0535-lengthinpicometers, def:physics:phyx-mini-0535:phyxminiproblems-problemphyxmini0535-answerchoice}
  A displayed answer is uniquely closest to a derived physical separation
\end{definition}
\begin{proof}
  This is the declared model definition of Is Unique Closest Displayed Separation.
\end{proof}
\begin{definition}[stationary Separation In Meters]
  \label{def:physics:phyx-mini-0535:phyxminiproblems-problemphyxmini0535-stationaryseparationinmeters}
  \lean{PhyXMiniProblems.ProblemPhyXMini0535.stationarySeparationInMeters}
  \uses{def:physics:phyx-mini-0535:phyxminiproblems-problemphyxmini0535-massinkilograms, def:physics:phyx-mini-0535:phyxminiproblems-problemphyxmini0535-chargeincoulombs, def:physics:phyx-mini-0535:phyxminiproblems-problemphyxmini0535-actioninjouleseconds, def:physics:phyx-mini-0535:phyxminiproblems-problemphyxmini0535-fourparticlecrystalsetup}
  The stationary positive length obtained by balancing the inverse-square rigid-box kinetic term with the attractive inverse-distance Coulomb term. This is an analytic candidate built from independent physical constants; its minimality and numerical answer-choice consequences are proved below rather than built into the definition
\end{definition}
\begin{proof}
  This is the declared model definition of stationary Separation In Meters.
\end{proof}
\begin{lemma}[total Energy Curve eq]
  \label{lem:physics:phyx-mini-0535:phyxminiproblems-problemphyxmini0535-totalenergycurve-eq}
  \lean{PhyXMiniProblems.ProblemPhyXMini0535.totalEnergyCurve_eq}
  \uses{def:physics:phyx-mini-0535:phyxminiproblems-problemphyxmini0535-lengthquantity, def:physics:phyx-mini-0535:phyxminiproblems-problemphyxmini0535-lengthinmeters, def:physics:phyx-mini-0535:phyxminiproblems-problemphyxmini0535-massinkilograms, def:physics:phyx-mini-0535:phyxminiproblems-problemphyxmini0535-chargeincoulombs, def:physics:phyx-mini-0535:phyxminiproblems-problemphyxmini0535-actioninjouleseconds, def:physics:phyx-mini-0535:phyxminiproblems-problemphyxmini0535-fourparticlecrystalsetup, def:physics:phyx-mini-0535:phyxminiproblems-problemphyxmini0535-matchestwoiontwoelectronscenario, def:physics:phyx-mini-0535:phyxminiproblems-problemphyxmini0535-matchesprimarycrystalfigure, def:physics:phyx-mini-0535:phyxminiproblems-problemphyxmini0535-satisfiesrigidboxgroundstatelaw, def:physics:phyx-mini-0535:phyxminiproblems-problemphyxmini0535-satisfiescoulombandtotalenergylaws, def:physics:phyx-mini-0535:phyxminiproblems-problemphyxmini0535-totalenergyinjoulesatspacing}
  The two ground-state kinetic terms and six Coulomb pairs reduce to E(d) = pi\textasciicircum{}2 hbar\textasciicircum{}2/(9 m d\textasciicircum{}2) - 7 k e\textasciicircum{}2/(3 d). This is derived from the all-pairs laws and figure geometry, not assumed
\end{lemma}
\begin{proof}
  The conclusion follows from the governing relations and figure readouts named in the statement of total Energy Curve eq.
\end{proof}
\begin{lemma}[exists unique minimizing separation]
  \label{lem:physics:phyx-mini-0535:phyxminiproblems-problemphyxmini0535-exists-unique-minimizing-separation}
  \lean{PhyXMiniProblems.ProblemPhyXMini0535.exists_unique_minimizing_separation}
  \uses{def:physics:phyx-mini-0535:phyxminiproblems-problemphyxmini0535-lengthquantity, def:physics:phyx-mini-0535:phyxminiproblems-problemphyxmini0535-lengthinmeters, def:physics:phyx-mini-0535:phyxminiproblems-problemphyxmini0535-fourparticlecrystalsetup, def:physics:phyx-mini-0535:phyxminiproblems-problemphyxmini0535-matchestwoiontwoelectronscenario, def:physics:phyx-mini-0535:phyxminiproblems-problemphyxmini0535-matchesprimarycrystalfigure, def:physics:phyx-mini-0535:phyxminiproblems-problemphyxmini0535-hasphysicalcrystalparameters, def:physics:phyx-mini-0535:phyxminiproblems-problemphyxmini0535-satisfiesrigidboxgroundstatelaw, def:physics:phyx-mini-0535:phyxminiproblems-problemphyxmini0535-satisfiescoulombandtotalenergylaws, def:physics:phyx-mini-0535:phyxminiproblems-problemphyxmini0535-isuniqueenergyminimizingseparation, def:physics:phyx-mini-0535:phyxminiproblems-problemphyxmini0535-stationaryseparationinmeters}
  For positive constants the curve above has its unique positive minimum at 2*pi\textasciicircum{}2*hbar\textasciicircum{}2/(21*k*m*e\textasciicircum{}2). This lemma does not use any displayed answer
\end{lemma}
\begin{proof}
  The conclusion follows from the governing relations and figure readouts named in the statement of exists unique minimizing separation.
\end{proof}
% --- Archon declaration topology end ---
% --- Archon physics formalization source end ---
